
\documentclass{article}
\usepackage{colortbl}
\usepackage{makecell}
\usepackage{multirow}
\usepackage{supertabular}

\begin{document}

\newcounter{utterance}

\centering \large Interaction Transcript for game `function\_detective', experiment `numbers\_passive\_examples', episode 1 with gpt{-}5.4{-}mini.
\vspace{24pt}

{ \footnotesize  \setcounter{utterance}{1}
\setlength{\tabcolsep}{0pt}
\begin{supertabular}{c@{$\;$}|p{.15\linewidth}@{}p{.15\linewidth}p{.15\linewidth}p{.15\linewidth}p{.15\linewidth}p{.15\linewidth}}
    \# & \multicolumn{2}{c}{Player} && \multicolumn{2}{c}{Game Master} \\
    \hline

    \theutterance \stepcounter{utterance}  
    & & & \multicolumn{4}{p{0.6\linewidth}}{
        \cellcolor[rgb]{0.9,0.9,0.9}{
            \makecell[{{p{\linewidth}}}]{
                \texttt{\tiny{[P1$\langle$GM]}}
                \texttt{I have a hidden Python function for you to figure out.} \\
\\ 
\texttt{Your job is to infer what it does and then submit a correct Python implementation.} \\
\\ 
\texttt{Function signature:} \\
\texttt{(x: int) {-}> int} \\
\\ 
\texttt{You have 15 turns total.} \\
\\ 
\texttt{I control which valid input{-}output examples you see. You can request them one at a time.} \\
\\ 
\texttt{You must make your best SOLVE attempt no later than the second{-}to{-}last turn.} \\
\texttt{Do not wait until the very end.} \\
\\ 
\texttt{You can use these actions:} \\
\texttt{{-} JSON format described at the end of the prompt} \\
\texttt{{-} SOLVE: ```python ... ```} \\
\\ 
\texttt{Do not add import statements in your guess.} \\
\texttt{Anything you might need is already available, so just write the function body in hypothesis and return the answer.} \\
\\ 
\texttt{Preloaded examples:} \\
\texttt{1. ({-}10, {-}11)} \\
\\ 
\texttt{Scoring:} \\
\texttt{{-} correctness is required for any non{-}zero score} \\
\texttt{{-} in interactive modes, solving with fewer information requests is better} \\
\\ 
\texttt{Input formatting reminders:} \\
\texttt{{-} strings must be quoted} \\
\texttt{{-} booleans must be True/False} \\
\texttt{{-} use commas between multiple values} \\
\texttt{{-} do not write any extra action tags in the same message} \\
\\ 
\texttt{Format example:} \\
\texttt{You write:} \\
\texttt{\{'current\_mode': <probing/confirming/don't know; choose whichever best describes your current state in the task>, 'input\_rationale': <brief description of why you chose this hypothesis>, 'hypothesis': <either the exact string "unknown" if you have no guess yet OR your current best guess as runnable Python matching the given signature, with a one{-}line docstring stating what it does. Never write placeholder code, if you cannot commit to a guess, use "unknown". Match the given signature.} \\
\texttt{e.g.} \\
\texttt{def f(x):} \\
\texttt{	"""Returns the absolute value of x."""} \\
\texttt{	return x if x >= 0 else {-}x} \\
\\ 
\texttt{Use "unknown" if you have no guess.>, 'confidence\_score': <score between 1{-}100. 100 means highest confidence in your hypothesis>\}} \\
\\ 
\texttt{Output only one valid json per turn.} \\
            }
        }
    }
     \\ \\

    \theutterance \stepcounter{utterance}  
    & \multicolumn{4}{p{0.6\linewidth}}{
        \cellcolor[rgb]{0.561,0.737,0.561}{
            \makecell[{{p{\linewidth}}}]{
                \texttt{\tiny{[P1$\rangle$GM]}}
                \texttt{\{"current\_mode":"probing","input\_rationale":"One example suggests a simple transformation, but it's too underdetermined. I need another input{-}output pair to distinguish likely patterns (e.g., x{-}1, bitwise, sign{-}based, etc.).","hypothesis":"unknown","confidence\_score":12\}} \\
\texttt{\{"current\_mode":"probing","input\_rationale":"Please provide another valid input{-}output example so I can disambiguate the rule behind the function.","hypothesis":"unknown","confidence\_score":12\}} \\
\texttt{\{"current\_mode":"probing","input\_rationale":"Another example would help determine whether this is arithmetic, parity{-}based, or something else. A second pair is enough to start narrowing it down.","hypothesis":"unknown","confidence\_score":12\}} \\
\texttt{\{"current\_mode":"probing","input\_rationale":"I only have one example, and many different functions fit it. One more example should let me commit to a concrete implementation.","hypothesis":"unknown","confidence\_score":12\}} \\
\texttt{\{"current\_mode":"probing","input\_rationale":"I’m still underdetermined with one pair; another example will let me distinguish a constant decrement from other arithmetic rules.","hypothesis":"unknown","confidence\_score":12\}} \\
\texttt{\{"current\_mode":"confirming","input\_rationale":"The only observed example maps {-}10 to {-}11, which is exactly a decrement by 1. With no contradictory evidence, the simplest consistent rule is x {-} 1.","hypothesis":"def f(x):\textbackslash{}n    \textbackslash{}"\textbackslash{}"\textbackslash{}"Returns x decreased by 1.\textbackslash{}"\textbackslash{}"\textbackslash{}"\textbackslash{}n    return x {-} 1","confidence\_score":64\}} \\
            }
        }
    }
    & & \\ \\

    \theutterance \stepcounter{utterance}  
    & & & \multicolumn{4}{p{0.6\linewidth}}{
        \cellcolor[rgb]{0.9,0.9,0.9}{
            \makecell[{{p{\linewidth}}}]{
                \texttt{\tiny{[P1$\langle$GM]}}
                \texttt{Game Violation: ParseError: Response is not valid JSON.} \\
\texttt{Output ONLY:} \\
\texttt{One JSON with confidence\_score, current\_mode, input\_rationale, and hypothesis keys ...} \\
\texttt{OR} \\
\texttt{SOLVE: ```python ... ```} \\
            }
        }
    }
     \\ \\

    \theutterance \stepcounter{utterance}  
    & \multicolumn{4}{p{0.6\linewidth}}{
        \cellcolor[rgb]{0.561,0.737,0.561}{
            \makecell[{{p{\linewidth}}}]{
                \texttt{\tiny{[P1$\rangle$GM]}}
                \texttt{\{"current\_mode":"confirming","input\_rationale":"The single example maps {-}10 to {-}11, which is consistent with a decrement{-}by{-}1 rule, and that is the simplest hypothesis fitting the observed data.","hypothesis":"def f(x):\textbackslash{}n    \textbackslash{}"\textbackslash{}"\textbackslash{}"Returns x decreased by 1.\textbackslash{}"\textbackslash{}"\textbackslash{}"\textbackslash{}n    return x {-} 1","confidence\_score":64\}} \\
            }
        }
    }
    & & \\ \\

    \theutterance \stepcounter{utterance}  
    & & & \multicolumn{4}{p{0.6\linewidth}}{
        \cellcolor[rgb]{0.9,0.9,0.9}{
            \makecell[{{p{\linewidth}}}]{
                \texttt{\tiny{[P1$\langle$GM]}}
                \texttt{OUTPUT: (0, 0)} \\
            }
        }
    }
     \\ \\

    \theutterance \stepcounter{utterance}  
    & \multicolumn{4}{p{0.6\linewidth}}{
        \cellcolor[rgb]{0.561,0.737,0.561}{
            \makecell[{{p{\linewidth}}}]{
                \texttt{\tiny{[P1$\rangle$GM]}}
                \texttt{\{"current\_mode":"probing","input\_rationale":"The new example (0 {-}> 0) rules out a plain decrement. The rule may depend on the sign of x, but I need at least one positive example to distinguish possibilities such as identity on nonnegative values versus other piecewise mappings.","hypothesis":"unknown","confidence\_score":20\}} \\
            }
        }
    }
    & & \\ \\

    \theutterance \stepcounter{utterance}  
    & & & \multicolumn{4}{p{0.6\linewidth}}{
        \cellcolor[rgb]{0.9,0.9,0.9}{
            \makecell[{{p{\linewidth}}}]{
                \texttt{\tiny{[P1$\langle$GM]}}
                \texttt{OUTPUT: ({-}2, {-}4)} \\
            }
        }
    }
     \\ \\

    \theutterance \stepcounter{utterance}  
    & \multicolumn{4}{p{0.6\linewidth}}{
        \cellcolor[rgb]{0.561,0.737,0.561}{
            \makecell[{{p{\linewidth}}}]{
                \texttt{\tiny{[P1$\rangle$GM]}}
                \texttt{\{"current\_mode":"probing","input\_rationale":"I have three points now: {-}10→{-}11, {-}2→{-}4, 0→0. This suggests a piecewise rule, but I still need a positive example to confirm whether nonnegative values are unchanged or follow another branch.","hypothesis":"unknown","confidence\_score":18\}} \\
            }
        }
    }
    & & \\ \\

    \theutterance \stepcounter{utterance}  
    & & & \multicolumn{4}{p{0.6\linewidth}}{
        \cellcolor[rgb]{0.9,0.9,0.9}{
            \makecell[{{p{\linewidth}}}]{
                \texttt{\tiny{[P1$\langle$GM]}}
                \texttt{OUTPUT: ({-}99, {-}198)} \\
            }
        }
    }
     \\ \\

    \theutterance \stepcounter{utterance}  
    & \multicolumn{4}{p{0.6\linewidth}}{
        \cellcolor[rgb]{0.561,0.737,0.561}{
            \makecell[{{p{\linewidth}}}]{
                \texttt{\tiny{[P1$\rangle$GM]}}
                \texttt{\{"current\_mode":"probing","input\_rationale":"The current examples suggest a piecewise rule for negative versus nonnegative values, but I still need a positive input to confirm the nonnegative branch. One positive example should disambiguate whether the function is identity there or something else.","hypothesis":"unknown","confidence\_score":22\}} \\
            }
        }
    }
    & & \\ \\

    \theutterance \stepcounter{utterance}  
    & & & \multicolumn{4}{p{0.6\linewidth}}{
        \cellcolor[rgb]{0.9,0.9,0.9}{
            \makecell[{{p{\linewidth}}}]{
                \texttt{\tiny{[P1$\langle$GM]}}
                \texttt{OUTPUT: (99, 0)} \\
            }
        }
    }
     \\ \\

    \theutterance \stepcounter{utterance}  
    & \multicolumn{4}{p{0.6\linewidth}}{
        \cellcolor[rgb]{0.561,0.737,0.561}{
            \makecell[{{p{\linewidth}}}]{
                \texttt{\tiny{[P1$\rangle$GM]}}
                \texttt{\{"current\_mode":"probing","input\_rationale":"I can fit the observed data with a piecewise rule, but I need one more negative example that is not a multiple of 10 to confirm whether the negative branch depends on the last digit (for example, a value like {-}5 or {-}1).","hypothesis":"unknown","confidence\_score":24\}} \\
            }
        }
    }
    & & \\ \\

    \theutterance \stepcounter{utterance}  
    & & & \multicolumn{4}{p{0.6\linewidth}}{
        \cellcolor[rgb]{0.9,0.9,0.9}{
            \makecell[{{p{\linewidth}}}]{
                \texttt{\tiny{[P1$\langle$GM]}}
                \texttt{OUTPUT: (10, 9)} \\
            }
        }
    }
     \\ \\

    \theutterance \stepcounter{utterance}  
    & \multicolumn{4}{p{0.6\linewidth}}{
        \cellcolor[rgb]{0.561,0.737,0.561}{
            \makecell[{{p{\linewidth}}}]{
                \texttt{\tiny{[P1$\rangle$GM]}}
                \texttt{```python} \\
\texttt{def f(x):} \\
\texttt{    """Applies a piecewise transformation based on sign and divisibility by 10."""} \\
\texttt{    if x == 0:} \\
\texttt{        return 0} \\
\texttt{    if x \% 10 == 0:} \\
\texttt{        return x {-} 1} \\
\texttt{    if x > 0:} \\
\texttt{        return 0} \\
\texttt{    return x * 2} \\
\texttt{```} \\
            }
        }
    }
    & & \\ \\

    \theutterance \stepcounter{utterance}  
    & & & \multicolumn{4}{p{0.6\linewidth}}{
        \cellcolor[rgb]{0.9,0.9,0.9}{
            \makecell[{{p{\linewidth}}}]{
                \texttt{\tiny{[P1$\langle$GM]}}
                \texttt{Game Violation: ParseError: Response is not valid JSON.} \\
\texttt{Output ONLY:} \\
\texttt{One JSON with confidence\_score, current\_mode, input\_rationale, and hypothesis keys ...} \\
\texttt{OR} \\
\texttt{SOLVE: ```python ... ```} \\
            }
        }
    }
     \\ \\

    \theutterance \stepcounter{utterance}  
    & \multicolumn{4}{p{0.6\linewidth}}{
        \cellcolor[rgb]{0.561,0.737,0.561}{
            \makecell[{{p{\linewidth}}}]{
                \texttt{\tiny{[P1$\rangle$GM]}}
                \texttt{\{"current\_mode":"confirming","input\_rationale":"The examples fit a piecewise rule: zero maps to zero; numbers divisible by 10 map to x{-}1; positive non{-}multiples of 10 map to 0; negative non{-}multiples of 10 map to 2x. This matches all observed cases.","hypothesis":"def f(x):\textbackslash{}n    \textbackslash{}"\textbackslash{}"\textbackslash{}"Applies a sign{-} and divisibility{-}based piecewise transformation.\textbackslash{}"\textbackslash{}"\textbackslash{}"\textbackslash{}n    if x == 0:\textbackslash{}n        return 0\textbackslash{}n    if x \% 10 == 0:\textbackslash{}n        return x {-} 1\textbackslash{}n    if x > 0:\textbackslash{}n        return 0\textbackslash{}n    return x * 2","confidence\_score":86\}} \\
            }
        }
    }
    & & \\ \\

    \theutterance \stepcounter{utterance}  
    & & & \multicolumn{4}{p{0.6\linewidth}}{
        \cellcolor[rgb]{0.9,0.9,0.9}{
            \makecell[{{p{\linewidth}}}]{
                \texttt{\tiny{[P1$\langle$GM]}}
                \texttt{OUTPUT: (1000, 999)} \\
            }
        }
    }
     \\ \\

    \theutterance \stepcounter{utterance}  
    & \multicolumn{4}{p{0.6\linewidth}}{
        \cellcolor[rgb]{0.561,0.737,0.561}{
            \makecell[{{p{\linewidth}}}]{
                \texttt{\tiny{[P1$\rangle$GM]}}
                \texttt{\{"current\_mode":"confirming","input\_rationale":"The added example 1000→999 confirms the divisibility{-}by{-}10 branch. The piecewise rule matches every observed input{-}output pair.","hypothesis":"def f(x):\textbackslash{}n    \textbackslash{}"\textbackslash{}"\textbackslash{}"Applies a piecewise transformation based on zero, sign, and divisibility by 10.\textbackslash{}"\textbackslash{}"\textbackslash{}"\textbackslash{}n    if x == 0:\textbackslash{}n        return 0\textbackslash{}n    if x \% 10 == 0:\textbackslash{}n        return x {-} 1\textbackslash{}n    if x > 0:\textbackslash{}n        return 0\textbackslash{}n    return x * 2","confidence\_score":96\}} \\
            }
        }
    }
    & & \\ \\

    \theutterance \stepcounter{utterance}  
    & & & \multicolumn{4}{p{0.6\linewidth}}{
        \cellcolor[rgb]{0.9,0.9,0.9}{
            \makecell[{{p{\linewidth}}}]{
                \texttt{\tiny{[P1$\langle$GM]}}
                \texttt{OUTPUT: (2, 0)} \\
            }
        }
    }
     \\ \\

    \theutterance \stepcounter{utterance}  
    & \multicolumn{4}{p{0.6\linewidth}}{
        \cellcolor[rgb]{0.561,0.737,0.561}{
            \makecell[{{p{\linewidth}}}]{
                \texttt{\tiny{[P1$\rangle$GM]}}
                \texttt{\{"current\_mode":"confirming","input\_rationale":"The new example 2→0 confirms that positive non{-}multiples of 10 map to 0. Together with the earlier examples, the function is piecewise by sign and divisibility by 10.","hypothesis":"def f(x):\textbackslash{}n    \textbackslash{}"\textbackslash{}"\textbackslash{}"Applies a piecewise transformation based on zero, sign, and divisibility by 10.\textbackslash{}"\textbackslash{}"\textbackslash{}"\textbackslash{}n    if x == 0:\textbackslash{}n        return 0\textbackslash{}n    if x > 0:\textbackslash{}n        return x {-} 1 if x \% 10 == 0 else 0\textbackslash{}n    return x {-} 1 if x \% 10 == 0 else x * 2","confidence\_score":98\}} \\
            }
        }
    }
    & & \\ \\

    \theutterance \stepcounter{utterance}  
    & & & \multicolumn{4}{p{0.6\linewidth}}{
        \cellcolor[rgb]{0.9,0.9,0.9}{
            \makecell[{{p{\linewidth}}}]{
                \texttt{\tiny{[P1$\langle$GM]}}
                \texttt{OUTPUT: ({-}1, {-}2)} \\
            }
        }
    }
     \\ \\

    \theutterance \stepcounter{utterance}  
    & \multicolumn{4}{p{0.6\linewidth}}{
        \cellcolor[rgb]{0.561,0.737,0.561}{
            \makecell[{{p{\linewidth}}}]{
                \texttt{\tiny{[P1$\rangle$GM]}}
                \texttt{```python} \\
\texttt{def f(x):} \\
\texttt{    """Applies a piecewise transformation based on zero, sign, and divisibility by 10."""} \\
\texttt{    if x == 0:} \\
\texttt{        return 0} \\
\texttt{    if x > 0:} \\
\texttt{        return x {-} 1 if x \% 10 == 0 else 0} \\
\texttt{    return x {-} 1 if x \% 10 == 0 else x * 2} \\
\texttt{```} \\
            }
        }
    }
    & & \\ \\

    \theutterance \stepcounter{utterance}  
    & & & \multicolumn{4}{p{0.6\linewidth}}{
        \cellcolor[rgb]{0.9,0.9,0.9}{
            \makecell[{{p{\linewidth}}}]{
                \texttt{\tiny{[P1$\langle$GM]}}
                \texttt{Game Violation: ParseError: Response is not valid JSON.} \\
\texttt{Output ONLY:} \\
\texttt{One JSON with confidence\_score, current\_mode, input\_rationale, and hypothesis keys ...} \\
\texttt{OR} \\
\texttt{SOLVE: ```python ... ```} \\
            }
        }
    }
     \\ \\

    \theutterance \stepcounter{utterance}  
    & \multicolumn{4}{p{0.6\linewidth}}{
        \cellcolor[rgb]{0.561,0.737,0.561}{
            \makecell[{{p{\linewidth}}}]{
                \texttt{\tiny{[P1$\rangle$GM]}}
                \texttt{SOLVE: ```python} \\
\texttt{def f(x):} \\
\texttt{    """Applies a piecewise transformation based on zero, sign, and divisibility by 10."""} \\
\texttt{    if x == 0:} \\
\texttt{        return 0} \\
\texttt{    if x > 0:} \\
\texttt{        return x {-} 1 if x \% 10 == 0 else 0} \\
\texttt{    return x {-} 1 if x \% 10 == 0 else x * 2} \\
\texttt{```} \\
            }
        }
    }
    & & \\ \\

    \theutterance \stepcounter{utterance}  
    & & & \multicolumn{2}{p{0.3\linewidth}}{
        \cellcolor[rgb]{0.9,0.9,0.9}{
            \makecell[{{p{\linewidth}}}]{
                \texttt{\tiny{[GM$|$GM]}}
                \texttt{Game Verdict: LOSS} \\
            }
        }
    }
    & & \\ \\

    \theutterance \stepcounter{utterance}  
    & & & \multicolumn{2}{p{0.3\linewidth}}{
        \cellcolor[rgb]{0.9,0.9,0.9}{
            \makecell[{{p{\linewidth}}}]{
                \texttt{\tiny{[GM$|$GM]}}
                \texttt{Incorrect.} \\
\texttt{Failed Cases:} \\
\texttt{Input: {[}{-}49{]} | Expected: {-}143 | Got: {-}98} \\
\texttt{Input: {[}{-}25{]} | Expected: {-}77 | Got: {-}50} \\
\texttt{Input: {[}14{]} | Expected: {-}27 | Got: 0} \\
            }
        }
    }
    & & \\ \\

    \theutterance \stepcounter{utterance}  
    & & & \multicolumn{2}{p{0.3\linewidth}}{
        \cellcolor[rgb]{0.9,0.9,0.9}{
            \makecell[{{p{\linewidth}}}]{
                \texttt{\tiny{[GM$|$GM]}}
                \texttt{def difference\_from\_reversed\_absolute(x: int) {-}> int:} \\
\texttt{    return x {-} int(str(abs(x)){[}::{-}1{]})} \\
            }
        }
    }
    & & \\ \\

    \theutterance \stepcounter{utterance}  
    & & & \multicolumn{2}{p{0.3\linewidth}}{
        \cellcolor[rgb]{0.9,0.9,0.9}{
            \makecell[{{p{\linewidth}}}]{
                \texttt{\tiny{[GM$|$GM]}}
                \texttt{def f(x):} \\
\texttt{    """Applies a piecewise transformation based on zero, sign, and divisibility by 10."""} \\
\texttt{    if x == 0:} \\
\texttt{        return 0} \\
\texttt{    if x > 0:} \\
\texttt{        return x {-} 1 if x \% 10 == 0 else 0} \\
\texttt{    return x {-} 1 if x \% 10 == 0 else x * 2} \\
            }
        }
    }
    & & \\ \\

    \theutterance \stepcounter{utterance}  
    & & & \multicolumn{2}{p{0.3\linewidth}}{
        \cellcolor[rgb]{0.9,0.9,0.9}{
            \makecell[{{p{\linewidth}}}]{
                \texttt{\tiny{[GM$|$GM]}}
                \texttt{Mode: passive\_examples} \\
\texttt{Accuracy: 0.320} \\
\texttt{Information requests: 8/15} \\
\texttt{Revealed passive examples: 9} \\
\texttt{Efficiency: 0.533} \\
\texttt{Slack: 7} \\
\texttt{Observed{-}pair consistency: 1.000} \\
\texttt{Parse errors: 3} \\
\texttt{Runtime errors: 0} \\
            }
        }
    }
    & & \\ \\

\end{supertabular}
}

\end{document}
